\chapter{Formal Underpinnings for Event Processing with Event Absence}
\label{ch:model}

This chapter formalizes the approach outlined in the previous chapter by introducing an abstract model of event-driven processing that explicitly accounts for missing observations during execution. The model defines the structure and semantics of events, characterizes expectations about observation timing, and specifies how event-driven computation proceeds when expected observations fail to arrive. To support continued reasoning under such conditions, the model incorporates mechanisms for reconstructing missing observations through state-based prediction and associating inferred events with confidence values that quantify the uncertainty introduced by inference.


\section{Event Model and Stream Semantics}
\label{sec:events-streams}

An \emph{event} represents an observation about an underlying system or process being monitored and is defined as
\begin{equation}
e = \langle id, type, src, t, v, conf, m \rangle ,
\end{equation}
\noindent
where $id$ uniquely identifies the event, $type$ represents the class of the event, $src$ identifies the source of the event, $t$ is the event time indicating when the observed phenomenon occurred, $v$ is the observed or inferred value, $conf \in \mathcal{C}$ is a confidence value associated with the event, and $m$ represents metadata.

An \emph{event stream} $E$ is a sequence of events
\begin{equation}
E = \langle e_1, e_2, \dots \rangle .
\end{equation}

Events originate from a set of \emph{event sources} $\mathcal{S}$, which correspond to external systems or processes being observed and produce values $v$ as observations over time. Each event source $s \in \mathcal{S}$ is characterized by an identifier, an event type, and an expected sampling interval:
\begin{equation}
s = \langle id_s, type_s, \Delta_s \rangle .
\end{equation}
\noindent
The expected sampling interval $\Delta_s$ specifies the time between successive observations.

Event sources act as \emph{event generators} by formatting observations as events and publishing them to the event stream $E$. For each event $e$ originating from source $s$, the source attribute is set to $e.src = id_s$, and the type attribute is set to $e.type = type_s$. Observed events are assigned a maximal confidence value according to the confidence scale $\mathcal{C}$.
Events published to the stream are made available to downstream \emph{event consumers}, which observe and react to events in $E$.


\section{Modeling Event Absence}
\label{sec:scheduling-missing}

According to \citet{wasserkrug2006taxonomy}, event streams are subject to occurrence uncertainty originating at their sources, such that observed events may represent only partial views of the system behaviour. In event-driven systems, this uncertainty must be resolved at runtime to determine whether an expected observation has occurred or whether its absence should be interpreted as a delay or as missing.

Let $\kappa$ be an \emph{event consumer} that has access to the expected sampling intervals $\Delta_s$ of event sources $s \in \mathcal{S}$. Using the current system time $t_{\mathrm{now}}$, $\kappa$ compares expected observation times with the appearance of events in the event stream to classify the status of expected observations.

For each source $s \in \mathcal{S}$, let $t^{\mathrm{last}}_s$ represent the event time of the most recent event in $E$ originating from $s$ that has been accounted for by $\kappa$. The next expected observation time for source $s$ is defined as
\begin{equation}
t^{\mathrm{exp}}_s = t^{\mathrm{last}}_s + \Delta_s .
\end{equation}
A tolerance window $\varepsilon_s$ is associated with source $s$ and defines the acceptable deviation from the expected observation time.

Each expected observation is classified according to the set
\begin{equation}
% \mathcal{O} = \{\textsc{Observed}, \textsc{Late}, \textsc{Missing}\}.
\mathcal{O} = \{\textsc{Observed}, \textsc{Missing}\}.
\end{equation}

An expected observation for source $s$ at time $t^{\mathrm{exp}}_s$ is classified as:
\begin{itemize}
    \item \textsc{Observed} if there exists an event $e \in E$ such that $e.src = id_s$, $e.type = type_s$, and $\lvert e.t - t^{\mathrm{exp}}_s \rvert \le \varepsilon_s$.
    % \item \textsc{Late} if there exists an event $e \in E$ such that $e.src = id_s$, $e.type = type_s$, and $\lvert e.t - t^{\mathrm{exp}}_s \rvert > \varepsilon_s$.
    \item \textsc{Missing} if $t_{\mathrm{now}} > t^{\mathrm{exp}}_s + \varepsilon_s$ and no such event has been observed.
\end{itemize}
\noindent




\section{Event Reconstruction and State-Based Inference}
\label{sec:state-based-inference}

For each event source $s \in \mathcal{S}$, let $r_s$ represent an \emph{event reconstruction mechanism} that acts as an \emph{event generator}. The behaviour of $r_s$ is governed by the classification of expected observations determined in \secref{sec:scheduling-missing}.
Each reconstructor $r_s$ contains a predictor with internal state $x_s$. As an event consumer, $r_s$ processes observed events associated with source $s$ in order to update this predictor state. As an event generator, $r_s$ emits reconstructed events corresponding to expected observations that are classified as missing.

The predictor supports two operations:
\begin{align}
\mathcal{U}_s\bigl(x_s^{i-1}, v_s^{i}\bigr)
&\quad\text{(update)}, \\
\hat{v}_s^{i} = \mathcal{P}_s\bigl(x_s^{i-1}\bigr)
&\quad\text{(prediction)}.
\end{align}

Each index $i$ corresponds to a successive expected observation for source $s$. When the expected observation is classified as \textsc{Observed}, the update operation $\mathcal{U}_s$ is applied, advancing the predictor state using the observed value $v_s^{i}$. When the expected observation is classified as \textsc{Missing}, the prediction operation $\mathcal{P}_s$ is applied, producing an estimated value $\hat{v}_s^{i}$ in the absence of an observation. Exactly one operation is applied at each step.

Uncertainty associated with inference is represented through a confidence value derived from the predictor state. A confidence function is defined as
\begin{equation}
conf_s^{i} = \mathcal{C}\bigl(x_s^{i}\bigr),
\end{equation}
\noindent
where confidence reflects the uncertainty accumulated by the predictor when operating without direct observation.

\section{Stream Augmentation}
\label{sec:confidence-stream-augmentation}

When an expected observation associated with source $s \in \mathcal{S}$ is classified as \textsc{Missing}, the reconstructor produces a replacement event defined as
\begin{equation}
\hat{e}_s^{i} =
\langle id, type_s, id_s, t^{\mathrm{exp}}_s, \hat{v}_s^{i}, conf_s^{i}, m_{\mathrm{pred}} \rangle ,
\end{equation}
\noindent
where $\hat{v}_s^{i}$ is the inferred value produced by the predictor, $conf_s^{i}$ is the associated confidence value derived from the predictor state, and $m_{\mathrm{pred}}$ records metadata indicating that the event was reconstructed rather than directly observed.

Reconstructed events are emitted into the event stream in the same manner as observed events. Observed and reconstructed events share a common structural representation, enabling uniform processing by downstream \emph{event consumers} and the incorporation of confidence-based queries.

The event stream is then extended to an augmented stream
\begin{equation}
E^{+} = E \cup \{ \hat{e}_s^{i} \},
\end{equation}




% ------------------------------------------
\section{Confidence-Aware Event Processing}
\label{sec:confidence-aware-semantics}

The augmented event stream $E^{+}$ produced by the model is intended to be consumed by downstream event-driven processing components. An \emph{event processor} is modeled as an \emph{event consumer} that observes events from $E^{+}$ and evaluates event patterns and temporal relations over the stream.

Let $e \in E^{+}$ be an event with an associated confidence value $conf(e) \in \mathcal{C}$. Event patterns and temporal relations may incorporate confidence values through additional predicates or aggregation functions, allowing processing logic to regulate pattern evaluation based on the confidence associated with inferred events.

At a basic level, confidence-aware processing may require individual events to satisfy a threshold condition
\begin{equation}
conf(e) \ge \theta ,
\end{equation}
\noindent
where $\theta$ is an application-defined confidence threshold. Such conditions allow event processors to suppress low-confidence events or to rank candidate matches according to their associated confidence values.

More complex patterns may combine the confidence of multiple events. Consider, for example, a temporal pattern requiring an event of type $A$ to be followed by an event of type $B$ within a time window $\tau$. Let $e_A$ and $e_B$ denote the corresponding events, with associated confidence values $conf(e_A)$ and $conf(e_B)$. The confidence of the composite pattern may be approximated by an aggregation function, for example
\begin{equation}
conf(A \rightarrow B) = conf(e_A) \cdot conf(e_B).
\end{equation}
Alternative aggregation strategies may also be used, including minimum, maximum, average, or recency-based confidence, depending on application requirements.

These mechanisms enable uncertainty introduced by reconstructed events to be accounted for explicitly during pattern evaluation, while preserving the structure of event-driven processing and requiring no changes to the event processing model. 

